\section{SEANCE 3}
L'objectif de cette séance est de récupérer le code d'origine du robot, le compiler, le téléverser, et vérifier que son
fonctionnement d'usine est correct.
On doit qualifier les différentes fonctions du code par leur contribution spécifique.
Les fonctions qui ne sont pas utilisées dans le cadre de l'objectif de la SAE doivent être identifiées et supprimées,
tout comme les valeurs numériques non nommées, les duplications et les chemins morts.
Un plan d'épuration doit être rédigé indiquant ce qui doit rester, ce qui doit être supprimé, et dans quel ordre.
Des tests de non-régression doivent être effectués après chaque étape.

Documents à fournir
tableau fonction rôle status (garder/supprimer), et plan d'épuration argumenté


Tumbller.ino

//////////////////////////////////
//////////////////////////////////
// IMPORTS
//////////////////////////////////

//////////////////////////////////
//////////////////////////////////
// VARIABLES GLOBALES / CONSTANTES
//////////////////////////////////

function_mode
motion_mode


deux variables d'états sont utilisées par la boucle 
principale d'execution. Plusieurs modes d'execution
sont prévus initialement.
Il prévoit la possibilité d'utiliser pour contrôler
les mouvements du robot, soit une  télécommande 
infrarouge, soit un émétteur bluetooth.

cette fonction va réaliser un branchement switch/case 
sur la variable function_mode qui est une enum (variable définissant plusieurs valeurs
différentiables les unes d'entre elles, mais rattachées à cette enum, réaliser une 
switch/capteurs sur une variable enum signifie donner une instruction case pour 
toutes les valeurs possibles de l'enum ainsi que le default).
En fonction de cette enum, on va executer le code correspondant à un état particulier du robot
(couché, équilibre, balancier, transition)
void functionMode();

cette fonction réalise deux opérations switch/case une après l'autre
et pour chaque cas, assigne une valeur de vitesse pour les variables 
de la consigne de vitesse linéaire, et vitesse de rotation.
Le premier switch est fait par rapport au type de déplacement, le deuxième est
fait par rapport au type de contrôle.
Pour le mode de mouvement STOP et START, l'état d'activation des leds sont également modifiés.
Pour le mode START, la fonction carStop est appelée pour couper la consigne de vitesse du robot
setMotionState()

cette fonction va réaliser un switch/case sur la variable caractère key_value, suivant la valeur
de cette variable, le robot va être bloqué dans le mode de mouvement correspondant. La variable
est remise à sa valeur de départ à la fin de la fonction, après le switch/case.
Pour maintenir le robot dans le mode correspondant, la variable devra être remise à la bonne valeur
à chaque execution de la boucle principale
keyEventHandle()

//////////////////////////////////////////////
//////////////////////////////////////////////
// fonction de la boucle d'execution d'arduino
//////////////////////////////////////////////

setup s'execute une seule fois au début de l'éxecution du programme.
Son fonctionnenment est le suivant : 
Serial.begin(9600) -> initialisation du port série 
ultrasonicInit() -> initialisation des pins et
interuptions du capteur ultrasonique 
keyInit() -> initialisation du pin et de l'interuption de la touche
permettant de lancer une routine de test des différents modes
 qui change toutes les 500ms
rgb.initialize() -> on initialise la led rgb programmable 
voltageInit() -> on configure la tension utilisée comme référence dans les
entrées analogiques
start_prev_time = millis() -> assignation du temps depuis l'execution du programme,
cette valeur sert pour mesurer des différences de temps et donc d'executer
du code à un temps donné
carInitialize() -> on initialise les pins restants et on assigne balanceCar comme 
fonction pour un objet timer. On stoppe également le robot dans  cette partie
setup()

loop s'execute continuellement
La boucle principale d'execution s'articule autour de la lecture de variables et grandeurs physiques
et valeurs de capteurs, le calcul d'une consigne cible, puis le changement d'état du robot en fonction
de son environment (présence obstacle, tourner).
getKeyValue() -> assigne une valeur à la variable key_value
getBluetoothData() -> assigne la variable function_mode (soit BLUETOOTH soit IDLE)
keyEventHandle() -> assigne des variables et notamment motion_mode, en fonction de key_value
getDistance() -> on effectue une condition sur une différence de temps (50ms),
donc toute les 50ms, on va effectuer ce qui est dans la condition
(raz du compteur, assignation de l'interrupt sur ECHO_PIN, et on effectue 3 
écritures, bas, haut, puis nouveau bas)
voltageMeasure() -> effectue une condition sur le niveau de la variable "voltage", alumage de la led
si le niveau est trop bas
setMotionState() -> assigne les consignes de vitesse linéaire et de rotation
functionMode() -> assigne la boucle principale, la routine, que le robot
va effectuer, c'est les fonctions appellées dans cette fonction là qui vont 
faire que le robot avance tout seul, ou essaie de suivre
checkObstacle() -> assigne les variables booléennes de présence d'obstacles, 
utilisées dans les fonctions des modes de fonctionnement du robot
rgb.blink(100) -> fait clignoter la led
bloc direct -> affichage de la valeur de l'angle du filtre kalman
bloc direct -> stop le robot
bloc direct -> passage en mode balancier du robot
loop()

Tumbller.ino
contient la boucle d'execution que la carte arduino va charger. Elle charge
aussi les bibliothèques principales et importe la logique utilisée pour
les calculs de vitesse

faire remonter les CONSTANTES
exposer plus facilement les paramètres
mettre de la poo pour faciliter la réusabilité, l'extension
encapsuler
meilleure séparation des responsabilitées
le flux d'information est difficilement lisible
mettre des interfaces